% GENERATED FILE. Edit canonical lesson frontmatter, then run npm run generate:books.
% canonical-entries: 98

\chapter*{Glossary}
\addcontentsline{toc}{chapter}{Glossary}
\markboth{Glossary}{Glossary}

This glossary collects every word and phrase taught in the book. Pronunciation appears when it differs from the written form; chapter numbers show where each entry is introduced.

\noindent\begin{minipage}[t]{\linewidth}
\raggedright
\textbf{accipiō, obtineō}\par
\small \textquotedblleft{}I get\textquotedblright{} — two verbs, each one a prefix bolted onto a verb you already know, and each one bent out of shape by the join\par
\footnotesize Introduced in Chapter 42.
\end{minipage}\par\medskip

\noindent\begin{minipage}[t]{\linewidth}
\raggedright
\textbf{adiuvō, adiuvāre}\par
\small \textquotedblleft{}I help\textquotedblright{} — the verb hiding inside English aid, aide and adjutant, and no relation to juvenile\par
\footnotesize Introduced in Chapter 39.
\end{minipage}\par\medskip

\noindent\begin{minipage}[t]{\linewidth}
\raggedright
\textbf{ambulō, ambulāre}\par
\small \textquotedblleft{}I walk\textquotedblright{} — the verb behind ambulance, preamble and the pram, and a rare case where the front half is clear and the back half is not\par
\footnotesize Introduced in Chapter 41.
\end{minipage}\par\medskip

\noindent\begin{minipage}[t]{\linewidth}
\raggedright
\textbf{amō, amāre}\par
\small \textquotedblleft{}I love\textquotedblright{} — the chapter's fourth verb, the first thing generations of pupils ever chanted, and the word inside enemy\par
\footnotesize Introduced in Chapter 39.
\end{minipage}\par\medskip

\noindent\begin{minipage}[t]{\linewidth}
\raggedright
\textbf{aperiō, aperīre}\par
\small \textquotedblleft{}I open\textquotedblright{} — half of a matched pair, built from one root with two opposite prefixes\par
\footnotesize Introduced in Chapter 41.
\end{minipage}\par\medskip

\noindent\begin{minipage}[t]{\linewidth}
\raggedright
\textbf{aqua, vīnum}\par
\small water and wine — one word French wore down to almost nothing, one that barely changed anywhere\par
\footnotesize Introduced in Chapter 10.
\end{minipage}\par\medskip

\noindent\begin{minipage}[t]{\linewidth}
\raggedright
\textbf{audiō, audīre}\par
\small \textquotedblleft{}I hear\textquotedblright{} — the verb inside audio, audience and audit, and, less obviously, inside obey\par
\footnotesize Introduced in Chapter 40.
\end{minipage}\par\medskip

\noindent\begin{minipage}[t]{\linewidth}
\raggedright
\textbf{avē / avēte}\par
\small hail (a formal greeting)\par
\footnotesize Introduced in Chapter 1.
\end{minipage}\par\medskip

\noindent\begin{minipage}[t]{\linewidth}
\raggedright
\textbf{bene}\par
\small \textquotedblleft{}well\textquotedblright{} — a morphologically simple inheritance: French bien and Spanish bien are just bene plus one regular sound change (the same short-e-diphthongizes-to-ie shift as mel$\to$miel), with no extra grafted-on suffix, unlike quōmodo's more roundabout Romance history\par
\footnotesize Introduced in Chapter 27.
\end{minipage}\par\medskip

\noindent\begin{minipage}[t]{\linewidth}
\raggedright
\textbf{bonam noctem}\par
\small \textquotedblleft{}good night\textquotedblright{} — genuine Latin grammar (wishes take the accusative, not the nominative), though the fixed phrase itself may be a later/modern convention rather than a directly-attested ancient formula\par
\footnotesize Introduced in Chapter 30.
\end{minipage}\par\medskip

\noindent\begin{minipage}[t]{\linewidth}
\raggedright
\textbf{bonum māne}\par
\small \textquotedblleft{}good morning\textquotedblright{} — a modern, pedagogical phrase (not classically attested); unlike bonam noctem, its grammar is genuinely invisible, since māne is neuter and Latin neuters never distinguish nominative from accusative — Romans instead had a real, richly documented morning GREETING RITUAL, the salūtātiō, whose root is a doublet of salvē's own\par
\footnotesize Introduced in Chapter 32.
\end{minipage}\par\medskip

\noindent\begin{minipage}[t]{\linewidth}
\raggedright
\textbf{bonum vesperum}\par
\small \textquotedblleft{}good evening\textquotedblright{} — a modern, pedagogical phrase (not classically attested), but unlike bonum māne, this one CAN show the real Accusative-of-Exclamation grammar again, since vesper is masculine, not neuter\par
\footnotesize Introduced in Chapter 34.
\end{minipage}\par\medskip

\noindent\begin{minipage}[t]{\linewidth}
\raggedright
\textbf{caelum, caelī}\par
\small the sky, and the heavens — the word under celestial, and under a ceiling painted like the sky\par
\footnotesize Introduced in Chapter 45.
\end{minipage}\par\medskip

\noindent\begin{minipage}[t]{\linewidth}
\raggedright
\textbf{canis, fēlēs, cattus}\par
\small dog and (two words for) cat — canis is the \textquotedblleft{}expected\textquotedblright{} ancestor Spanish's perro mysteriously replaced; fēlēs was Classical Latin's OWN word for cat, later pushed out by the newcomer cattus\par
\footnotesize Introduced in Chapter 17.
\end{minipage}\par\medskip

\noindent\begin{minipage}[t]{\linewidth}
\raggedright
\textbf{capiō, capere}\par
\small \textquotedblleft{}I take\textquotedblright{} — the verb the habeō lesson promised, and the true relative of English have\par
\footnotesize Introduced in Chapter 39.
\end{minipage}\par\medskip

\noindent\begin{minipage}[t]{\linewidth}
\raggedright
\textbf{caput}\par
\small the head — the word French abandoned, but Spanish and English kept working hard\par
\footnotesize Introduced in Chapter 8.
\end{minipage}\par\medskip

\noindent\begin{minipage}[t]{\linewidth}
\raggedright
\textbf{claudō, claudere}\par
\small \textquotedblleft{}I close\textquotedblright{} — the word English calls \emph{close}, and the fourth of four\par
\footnotesize Introduced in Chapter 41.
\end{minipage}\par\medskip

\noindent\begin{minipage}[t]{\linewidth}
\raggedright
\textbf{cōgitō, cōgitāre}\par
\small I think" — the verb of \emph{cōgitō ergō sum}, built from a word meaning "to shake together\par
\footnotesize Introduced in Chapter 38.
\end{minipage}\par\medskip

\noindent\begin{minipage}[t]{\linewidth}
\raggedright
\textbf{conveniō, occurrō}\par
\small \textquotedblleft{}I meet\textquotedblright{} — two verbs for meeting, each one a prefix on a verb you already know, and the verb hiding inside a phrase from chapter 21\par
\footnotesize Introduced in Chapter 43.
\end{minipage}\par\medskip

\noindent\begin{minipage}[t]{\linewidth}
\raggedright
\textbf{crās tē vidēbō}\par
\small \textquotedblleft{}see you tomorrow\textquotedblright{} — built from two genuinely classical pieces (crās, \textquotedblleft{}tomorrow,\textquotedblright{} and vidēbō, already met last lesson) using the exact same pattern as the attested propediem tē vidēbō, though this specific combination isn't itself pulled from one surviving letter\par
\footnotesize Introduced in Chapter 25.
\end{minipage}\par\medskip

\noindent\begin{minipage}[t]{\linewidth}
\raggedright
\textbf{currō, currere}\par
\small \textquotedblleft{}I run\textquotedblright{} — the verb inside current, curriculum, courier and corridor, and, by way of a Celtic wagon, inside car\par
\footnotesize Introduced in Chapter 41.
\end{minipage}\par\medskip

\noindent\begin{minipage}[t]{\linewidth}
\raggedright
\textbf{dīcō, dīcere}\par
\small \textquotedblleft{}I say\textquotedblright{} — a verb that originally meant \textquotedblleft{}point out\textquotedblright{}, which is why English got both dictionary and teach from it\par
\footnotesize Introduced in Chapter 37.
\end{minipage}\par\medskip

\noindent\begin{minipage}[t]{\linewidth}
\raggedright
\textbf{diēs}\par
\small \textquotedblleft{}day\textquotedblright{} — the word you've already been using all along, inside diēs Lūnae and medius diēs, now on its own\par
\footnotesize Introduced in Chapter 28.
\end{minipage}\par\medskip

\noindent\begin{minipage}[t]{\linewidth}
\raggedright
\textbf{diēs Lūnae, Mārtis, Mercuriī, Iovis, Veneris}\par
\small the weekdays (Monday–Friday) — the full, un-worn-down planet-god phrases\par
\footnotesize Introduced in Chapter 5.
\end{minipage}\par\medskip

\noindent\begin{minipage}[t]{\linewidth}
\raggedright
\textbf{diēs Saturnī, diēs Sōlis}\par
\small Saturday and Sunday — Rome's own pagan names, and how Christianity later renamed them\par
\footnotesize Introduced in Chapter 5.
\end{minipage}\par\medskip

\noindent\begin{minipage}[t]{\linewidth}
\raggedright
\textbf{dō, dare}\par
\small \textquotedblleft{}I give\textquotedblright{} — the last of the chapter's eight verbs, and the one that gave English data, date, tradition and dose\par
\footnotesize Introduced in Chapter 37.
\end{minipage}\par\medskip

\noindent\begin{minipage}[t]{\linewidth}
\raggedright
\textbf{domus, domūs}\par
\small the house — the building a family lives in, and the root under domestic, dome and dominate\par
\footnotesize Introduced in Chapter 44.
\end{minipage}\par\medskip

\noindent\begin{minipage}[t]{\linewidth}
\raggedright
\textbf{dormiō, dormīre}\par
\small \textquotedblleft{}I sleep\textquotedblright{} — the verb behind dormitory and dormant, and the one word English may only look as though it borrowed\par
\footnotesize Introduced in Chapter 40.
\end{minipage}\par\medskip

\noindent\begin{minipage}[t]{\linewidth}
\raggedright
\textbf{emō, emere}\par
\small \textquotedblleft{}I buy\textquotedblright{} — a verb that once meant simply \textquotedblleft{}I take\textquotedblright{}, which is the only reason redeem, example and premium mean what they do\par
\footnotesize Introduced in Chapter 43.
\end{minipage}\par\medskip

\noindent\begin{minipage}[t]{\linewidth}
\raggedright
\textbf{eō, īre}\par
\small \textquotedblleft{}I go\textquotedblright{} — the shortest verb in Latin, whose bare stem is a single vowel\par
\footnotesize Introduced in Chapter 37.
\end{minipage}\par\medskip

\noindent\begin{minipage}[t]{\linewidth}
\raggedright
\textbf{exspectō, exspectāre}\par
\small \textquotedblleft{}I wait\textquotedblright{} — to wait, in Latin, is to look out for something, and the four verbs of this chapter together\par
\footnotesize Introduced in Chapter 42.
\end{minipage}\par\medskip

\noindent\begin{minipage}[t]{\linewidth}
\raggedright
\textbf{ferō, ferre}\par
\small \textquotedblleft{}I bring, I carry\textquotedblright{} — one verb built out of two unrelated roots, and the parent of transfer, refer, prefer, differ, offer and suffer\par
\footnotesize Introduced in Chapter 42.
\end{minipage}\par\medskip

\noindent\begin{minipage}[t]{\linewidth}
\raggedright
\textbf{focus, focī}\par
\small the hearth — the fire in the middle of the house, and the reason a lens has a focus\par
\footnotesize Introduced in Chapter 44.
\end{minipage}\par\medskip

\noindent\begin{minipage}[t]{\linewidth}
\raggedright
\textbf{frater, soror}\par
\small brother and sister — the words Vulgar Latin later abandoned in favor of germanus\par
\footnotesize Introduced in Chapter 7.
\end{minipage}\par\medskip

\noindent\begin{minipage}[t]{\linewidth}
\raggedright
\textbf{grātiās (tibi) agō}\par
\small thank you (lit. \textquotedblleft{}I give thanks to you\textquotedblright{})\par
\footnotesize Introduced in Chapter 1.
\end{minipage}\par\medskip

\noindent\begin{minipage}[t]{\linewidth}
\raggedright
\textbf{habeō, habēre}\par
\small \textquotedblleft{}I have\textquotedblright{} — a verb that looks exactly like its English translation and is not related to it at all\par
\footnotesize Introduced in Chapter 37.
\end{minipage}\par\medskip

\noindent\begin{minipage}[t]{\linewidth}
\raggedright
\textbf{hōra}\par
\small hour — the source word itself, and the Roman hour was NOT a fixed length\par
\footnotesize Introduced in Chapter 13.
\end{minipage}\par\medskip

\noindent\begin{minipage}[t]{\linewidth}
\raggedright
\textbf{hortus, hortī}\par
\small the garden — an enclosed piece of ground, and a true blood relative of English yard and garden\par
\footnotesize Introduced in Chapter 44.
\end{minipage}\par\medskip

\noindent\begin{minipage}[t]{\linewidth}
\raggedright
\textbf{Iānuārius, Februārius, Mārtius, Aprīlis, Māius, Iūnius, Quīntīlis, Sextīlis, September, Octōber, November, December}\par
\small the twelve months — and the two ORIGINAL names, Quintilis and Sextilis, before Rome renamed them\par
\footnotesize Introduced in Chapter 11.
\end{minipage}\par\medskip

\noindent\begin{minipage}[t]{\linewidth}
\raggedright
\textbf{ignis, ignis}\par
\small fire itself — the classical word, which the Romance languages abandoned and English kept\par
\footnotesize Introduced in Chapter 45.
\end{minipage}\par\medskip

\noindent\begin{minipage}[t]{\linewidth}
\raggedright
\textbf{ignōsce mihi}\par
\small forgive me / I'm sorry (imperative of ignōscere, \textquotedblleft{}to not-know, overlook, pardon\textquotedblright{})\par
\footnotesize Introduced in Chapter 4.
\end{minipage}\par\medskip

\noindent\begin{minipage}[t]{\linewidth}
\raggedright
\textbf{intellegō, intellegere}\par
\small \textquotedblleft{}I understand\textquotedblright{} — literally \emph{inter} plus \emph{legō}, \textquotedblleft{}to read between\textquotedblright{}, and the ancestor of intellect and intelligence\par
\footnotesize Introduced in Chapter 38.
\end{minipage}\par\medskip

\noindent\begin{minipage}[t]{\linewidth}
\raggedright
\textbf{ita / nōn}\par
\small yes / no (lit. \textquotedblleft{}thus\textquotedblright{} / \textquotedblleft{}not\textquotedblright{})\par
\footnotesize Introduced in Chapter 1.
\end{minipage}\par\medskip

\noindent\begin{minipage}[t]{\linewidth}
\raggedright
\textbf{lac, lactis}\par
\small milk — the root of lactose, lettuce and, at one remove, of the Milky Way itself\par
\footnotesize Introduced in Chapter 47.
\end{minipage}\par\medskip

\noindent\begin{minipage}[t]{\linewidth}
\raggedright
\textbf{legō, legere}\par
\small \textquotedblleft{}I read\textquotedblright{} — a verb that first meant \textquotedblleft{}to gather\textquotedblright{}, which is why English got collect, elegant, lecture and legible from it\par
\footnotesize Introduced in Chapter 38.
\end{minipage}\par\medskip

\noindent\begin{minipage}[t]{\linewidth}
\raggedright
\textbf{lūdō, lūdere}\par
\small \textquotedblleft{}I play\textquotedblright{} — a small verb with an outsized English family, most of which is about being fooled rather than about play\par
\footnotesize Introduced in Chapter 42.
\end{minipage}\par\medskip

\noindent\begin{minipage}[t]{\linewidth}
\raggedright
\textbf{lūna, lūnae}\par
\small the moon — the same root as lux, worn down to mean 'the shining one'\par
\footnotesize Introduced in Chapter 46.
\end{minipage}\par\medskip

\noindent\begin{minipage}[t]{\linewidth}
\raggedright
\textbf{lūx, lūcis}\par
\small light — the root of lucid, translucent and Lucifer, and a cousin of English light itself\par
\footnotesize Introduced in Chapter 46.
\end{minipage}\par\medskip

\noindent\begin{minipage}[t]{\linewidth}
\raggedright
\textbf{māne}\par
\small \textquotedblleft{}morning, early\textquotedblright{} — the word already hiding inside dē māne (the source of French demain, Italian domani, and Catalan demà, all rebuilt \textquotedblleft{}tomorrow\textquotedblright{} from this phrase) and, without the dē-, the direct base of Spanish mañana too, now standing alone\par
\footnotesize Introduced in Chapter 31.
\end{minipage}\par\medskip

\noindent\begin{minipage}[t]{\linewidth}
\raggedright
\textbf{manus}\par
\small the hand — a rare feminine noun in an overwhelmingly masculine group of Latin words\par
\footnotesize Introduced in Chapter 8.
\end{minipage}\par\medskip

\noindent\begin{minipage}[t]{\linewidth}
\raggedright
\textbf{mare, maris}\par
\small the sea — the word that finishes Mediterranean and names every marine thing in English\par
\footnotesize Introduced in Chapter 45.
\end{minipage}\par\medskip

\noindent\begin{minipage}[t]{\linewidth}
\raggedright
\textbf{medius diēs, media nox}\par
\small noon and midnight — the source phrase daughters kept whole, wore down unevenly, or eroded almost entirely\par
\footnotesize Introduced in Chapter 12.
\end{minipage}\par\medskip

\noindent\begin{minipage}[t]{\linewidth}
\raggedright
\textbf{mel, mellis}\par
\small honey — the root of mellifluous and marmalade, and of molasses by way of Portuguese\par
\footnotesize Introduced in Chapter 47.
\end{minipage}\par\medskip

\noindent\begin{minipage}[t]{\linewidth}
\raggedright
\textbf{merīdiēs}\par
\small \textquotedblleft{}midday, noon\textquotedblright{} — the same medius diēs you already know, contracted into one word; Romans split the day around IT (ante/post merīdiem), not into a separate colloquial \textquotedblleft{}afternoon\textquotedblright{} the way Spanish's tarde or French's après-midi do\par
\footnotesize Introduced in Chapter 35.
\end{minipage}\par\medskip

\noindent\begin{minipage}[t]{\linewidth}
\raggedright
\textbf{mūrus, mūrī}\par
\small the wall — the one that keeps things out, not the one that divides a room\par
\footnotesize Introduced in Chapter 44.
\end{minipage}\par\medskip

\noindent\begin{minipage}[t]{\linewidth}
\raggedright
\textbf{niger, albus}\par
\small black and white — the source words for Spanish/French/Italian, each with a shine-vs-matte twin\par
\footnotesize Introduced in Chapter 6.
\end{minipage}\par\medskip

\noindent\begin{minipage}[t]{\linewidth}
\raggedright
\textbf{nihil est}\par
\small \textquotedblleft{}you're welcome\textquotedblright{} — another honest gap: Classical Latin has no single fixed reply to thanks, and nihil est (\textquotedblleft{}it's nothing\textquotedblright{}) is a modern conversational-Latin convention built from genuinely real words, not a specific classical citation for this exact job\par
\footnotesize Introduced in Chapter 23.
\end{minipage}\par\medskip

\noindent\begin{minipage}[t]{\linewidth}
\raggedright
\textbf{nox}\par
\small \textquotedblleft{}night\textquotedblright{} — the word already hiding inside media nox ('midnight'), one of the most rock-solid PIE roots in the whole Indo-European family\par
\footnotesize Introduced in Chapter 29.
\end{minipage}\par\medskip

\noindent\begin{minipage}[t]{\linewidth}
\raggedright
\textbf{ōvum, ōvī}\par
\small the egg — a cousin of English egg, with two sound changes that hide the relationship completely\par
\footnotesize Introduced in Chapter 47.
\end{minipage}\par\medskip

\noindent\begin{minipage}[t]{\linewidth}
\raggedright
\textbf{pānis}\par
\small bread — the source of \textquotedblleft{}companion,\textquotedblright{} and of the friendship-word every Romance daughter still shares with it\par
\footnotesize Introduced in Chapter 10.
\end{minipage}\par\medskip

\noindent\begin{minipage}[t]{\linewidth}
\raggedright
\textbf{pater, mater}\par
\small father and mother — the PIE ancestor Latin and English both inherited, in parallel\par
\footnotesize Introduced in Chapter 7.
\end{minipage}\par\medskip

\noindent\begin{minipage}[t]{\linewidth}
\raggedright
\textbf{piscis, piscis}\par
\small the fish — and the single clearest demonstration in this book of the sound law linking Latin to English\par
\footnotesize Introduced in Chapter 47.
\end{minipage}\par\medskip

\noindent\begin{minipage}[t]{\linewidth}
\raggedright
\textbf{porta, portae}\par
\small the door or gate — and the reason a harbour, a porch and a portal are all the same word\par
\footnotesize Introduced in Chapter 44.
\end{minipage}\par\medskip

\noindent\begin{minipage}[t]{\linewidth}
\raggedright
\textbf{propediem tē vidēbō}\par
\small \textquotedblleft{}see you later/soon\textquotedblright{} — a genuine change of pace from the last few lessons: this one IS solidly attested, twice, in Cicero's own letters\par
\footnotesize Introduced in Chapter 24.
\end{minipage}\par\medskip

\noindent\begin{minipage}[t]{\linewidth}
\raggedright
\textbf{quaesō}\par
\small please (literally \textquotedblleft{}I ask, I beg\textquotedblright{})\par
\footnotesize Introduced in Chapter 3.
\end{minipage}\par\medskip

\noindent\begin{minipage}[t]{\linewidth}
\raggedright
\textbf{quid agis?}\par
\small how are you? — literally, what are you doing?\par
\footnotesize Introduced in Chapter 19.
\end{minipage}\par\medskip

\noindent\begin{minipage}[t]{\linewidth}
\raggedright
\textbf{quid tibi nōmen est? / mihi nōmen est...}\par
\small what's your name? / my name is... — the genuinely classical dative of possession, literally what is the name to you?\par
\footnotesize Introduced in Chapter 20.
\end{minipage}\par\medskip

\noindent\begin{minipage}[t]{\linewidth}
\raggedright
\textbf{quōmodo}\par
\small how, in what way" — the genuine classical word behind Italian come, Spanish cómo, and French comment, which went one step further and double-marked "manner\par
\footnotesize Introduced in Chapter 26.
\end{minipage}\par\medskip

\noindent\begin{minipage}[t]{\linewidth}
\raggedright
\textbf{respondeō, respondēre}\par
\small \textquotedblleft{}I answer\textquotedblright{} — literally \textquotedblleft{}I pledge back\textquotedblright{}, from the verb Romans used for a solemn promise, and the reason a husband is a spouse\par
\footnotesize Introduced in Chapter 43.
\end{minipage}\par\medskip

\noindent\begin{minipage}[t]{\linewidth}
\raggedright
\textbf{rogō, rogāre}\par
\small \textquotedblleft{}I ask\textquotedblright{} — a verb probably built on \textquotedblleft{}stretching straight out\textquotedblright{}, and a cousin of regal, direct and right\par
\footnotesize Introduced in Chapter 39.
\end{minipage}\par\medskip

\noindent\begin{minipage}[t]{\linewidth}
\raggedright
\textbf{ruber, caeruleus}\par
\small red (rock-solid) and blue (surprisingly shaky) — Latin's clearest color word, and one of its weakest\par
\footnotesize Introduced in Chapter 6.
\end{minipage}\par\medskip

\noindent\begin{minipage}[t]{\linewidth}
\raggedright
\textbf{sāl, salis}\par
\small salt — a cousin of English salt, and the word a Roman soldier's salary was named for\par
\footnotesize Introduced in Chapter 47.
\end{minipage}\par\medskip

\noindent\begin{minipage}[t]{\linewidth}
\raggedright
\textbf{salvē (post merīdiem)}\par
\small why no two-word good-afternoon phrase exists — Latin has no simplex afternoon noun, and a grammatical three-word workaround is not attested as a greeting\par
\footnotesize Introduced in Chapter 36.
\end{minipage}\par\medskip

\noindent\begin{minipage}[t]{\linewidth}
\raggedright
\textbf{salvē / salvēte}\par
\small hello (lit. \textquotedblleft{}be well\textquotedblright{})\par
\footnotesize Introduced in Chapter 1.
\end{minipage}\par\medskip

\noindent\begin{minipage}[t]{\linewidth}
\raggedright
\textbf{sciō, scīre}\par
\small \textquotedblleft{}I know\textquotedblright{} — one of Latin's two verbs for knowing, and the one that gave English science and nice\par
\footnotesize Introduced in Chapter 37.
\end{minipage}\par\medskip

\noindent\begin{minipage}[t]{\linewidth}
\raggedright
\textbf{scrībō, scrībere}\par
\small \textquotedblleft{}I write\textquotedblright{} — a verb whose root means \textquotedblleft{}to scratch\textquotedblright{}, and the fourth of this chapter's four\par
\footnotesize Introduced in Chapter 38.
\end{minipage}\par\medskip

\noindent\begin{minipage}[t]{\linewidth}
\raggedright
\textbf{sedeō, sedēre}\par
\small \textquotedblleft{}I sit\textquotedblright{} — the verb that is English \emph{sit} by descent and English \emph{session} by loan, and the clearest case of the difference\par
\footnotesize Introduced in Chapter 40.
\end{minipage}\par\medskip

\noindent\begin{minipage}[t]{\linewidth}
\raggedright
\textbf{sex septem octō novem decem}\par
\small six to ten — and the four that got frozen into the calendar\par
\footnotesize Introduced in Chapter 2.
\end{minipage}\par\medskip

\noindent\begin{minipage}[t]{\linewidth}
\raggedright
\textbf{somnus, somnī}\par
\small sleep — the word under insomnia and somnolent, and a cousin of Sanskrit svapna\par
\footnotesize Introduced in Chapter 46.
\end{minipage}\par\medskip

\noindent\begin{minipage}[t]{\linewidth}
\raggedright
\textbf{stēlla, stēllae}\par
\small the star — a cousin of English star, and the word inside constellation and disaster\par
\footnotesize Introduced in Chapter 46.
\end{minipage}\par\medskip

\noindent\begin{minipage}[t]{\linewidth}
\raggedright
\textbf{stō, stāre}\par
\small \textquotedblleft{}I stand\textquotedblright{} — the shortest verb in the chapter and the most productive in English, and the fourth of four\par
\footnotesize Introduced in Chapter 40.
\end{minipage}\par\medskip

\noindent\begin{minipage}[t]{\linewidth}
\raggedright
\textbf{sum, esse}\par
\small \textquotedblleft{}I am\textquotedblright{} — the verb every sentence leans on, and the one whose forms English still shares outright\par
\footnotesize Introduced in Chapter 37.
\end{minipage}\par\medskip

\noindent\begin{minipage}[t]{\linewidth}
\raggedright
\textbf{tempestās}\par
\small weather or storm — a derivative of tempus that specialized for weather in Classical Latin, while tempus itself meant only time; later tempus widened to weather and tempestās narrowed to storm\par
\footnotesize Introduced in Chapter 15.
\end{minipage}\par\medskip

\noindent\begin{minipage}[t]{\linewidth}
\raggedright
\textbf{terra, terrae}\par
\small the earth — dry land as opposed to sea, and the root under terrain, terrace and Mediterranean\par
\footnotesize Introduced in Chapter 45.
\end{minipage}\par\medskip

\noindent\begin{minipage}[t]{\linewidth}
\raggedright
\textbf{tū / vōs}\par
\small you (singular) / you (plural) — a PURELY grammatical number distinction in Classical Latin, with no politeness dimension at all; the polite-vous pattern (a Late Latin innovation) is vōs's own direct descendant, unlike German's Sie, which solved the same politeness problem independently via a completely different pronoun (3rd-person plural, not 2nd)\par
\footnotesize Introduced in Chapter 22.
\end{minipage}\par\medskip

\noindent\begin{minipage}[t]{\linewidth}
\raggedright
\textbf{umbra, umbrae}\par
\small the shade or shadow — and the reason a little shade is called an umbrella\par
\footnotesize Introduced in Chapter 46.
\end{minipage}\par\medskip

\noindent\begin{minipage}[t]{\linewidth}
\raggedright
\textbf{ūndecim — vīgintī}\par
\small 11-20, the source numbers themselves — including Latin's own genuine oddity, SUBTRACTIVE 18 and 19, which Spanish quietly regularized away\par
\footnotesize Introduced in Chapter 16.
\end{minipage}\par\medskip

\noindent\begin{minipage}[t]{\linewidth}
\raggedright
\textbf{ūnus duo trēs quattuor quīnque}\par
\small one to five — the Roman originals that became every Romance number\par
\footnotesize Introduced in Chapter 2.
\end{minipage}\par\medskip

\noindent\begin{minipage}[t]{\linewidth}
\raggedright
\textbf{ut valēs?}\par
\small how do you fare?\par
\footnotesize Introduced in Chapter 19.
\end{minipage}\par\medskip

\noindent\begin{minipage}[t]{\linewidth}
\raggedright
\textbf{valē / valēte}\par
\small goodbye (lit. \textquotedblleft{}be strong\textquotedblright{})\par
\footnotesize Introduced in Chapter 1.
\end{minipage}\par\medskip

\noindent\begin{minipage}[t]{\linewidth}
\raggedright
\textbf{veniō, venīre}\par
\small \textquotedblleft{}I come\textquotedblright{} — a Latin verb that really is the same inherited word as English come\par
\footnotesize Introduced in Chapter 37.
\end{minipage}\par\medskip

\noindent\begin{minipage}[t]{\linewidth}
\raggedright
\textbf{ventus, ventī}\par
\small the wind — and a true Germanic cousin, with the sound change that proves it\par
\footnotesize Introduced in Chapter 45.
\end{minipage}\par\medskip

\noindent\begin{minipage}[t]{\linewidth}
\raggedright
\textbf{vēr, aestās, autumnus, hiems}\par
\small the four seasons — and why Romance languages often skipped the plain noun for an adjective phrase instead\par
\footnotesize Introduced in Chapter 9.
\end{minipage}\par\medskip

\noindent\begin{minipage}[t]{\linewidth}
\raggedright
\textbf{vesper}\par
\small evening — a genuine everyday classical word sharing a deep root with English west through the direction of sunset\par
\footnotesize Introduced in Chapter 33.
\end{minipage}\par\medskip

\noindent\begin{minipage}[t]{\linewidth}
\raggedright
\textbf{videō, vidēre}\par
\small \textquotedblleft{}I see\textquotedblright{} — the widest root in this chapter, because seeing and knowing were once one idea\par
\footnotesize Introduced in Chapter 37.
\end{minipage}\par\medskip

\noindent\begin{minipage}[t]{\linewidth}
\raggedright
\textbf{vīgintī annōs nātus sum}\par
\small \textquotedblleft{}I am twenty years old\textquotedblright{} — literally \textquotedblleft{}I am, having been born, twenty years {[}so far{]}\textquotedblright{}; NEITHER Romance's \textquotedblleft{}have\textquotedblright{} NOR Germanic's \textquotedblleft{}be my years\textquotedblright{} — Latin used a THIRD construction, later abandoned\par
\footnotesize Introduced in Chapter 14.
\end{minipage}\par\medskip

\noindent\begin{minipage}[t]{\linewidth}
\raggedright
\textbf{viridis, flāvus}\par
\small green and yellow — viridis is the direct, kept source of Spanish's verde; flāvus is Latin's own real word for yellow, and almost NO Romance language kept it as their everyday color word — most (French, Italian, Romanian) converged on a DIFFERENT shared Latin word instead\par
\footnotesize Introduced in Chapter 18.
\end{minipage}\par\medskip

\noindent\begin{minipage}[t]{\linewidth}
\raggedright
\textbf{volup est convēnisse}\par
\small \textquotedblleft{}it is a pleasure to have met you\textquotedblright{} — an attested Plautine phrase that can fill a conversational slot for which Classical Latin has no single fixed idiom\par
\footnotesize Introduced in Chapter 21.
\end{minipage}\par\medskip
